Fix \(P\in\mathfrak P_{\mathcal D}\) in the following causal subclass.  For every \(d\in\mathcal D\), the observed bridges \(h_d\) and \(q_d\) also satisfy the two latent bridge equations
\[
 \mathbb E_P[Y\mid S_{2d},A,U,X,G=O]
 =\mathbb E_P[h_d(S_{3d},S_{2d},A,X)\mid S_{2d},A,U,X,G=O]
\]
and
\[
 \frac{p_P(S_{2d},U,X\mid A,G=E)}{p_P(S_{2d},U,X\mid A,G=O)}
 =\mathbb E_P[q_d(S_{2d},S_{1d},A,X)\mid S_{2d},A,U,X,G=O].
\]
Consequently both Imbens--Kallus--Mao--Wang identification routes hold and \(\mu_{d,1}-\mu_{d,0}=\tau\) for every \(d\in\mathcal D\).

Use a fixed integer \(L\ge2\) of approximately equal observational folds.  Only the observational sample is split: \(\widehat h_{d,B}^{(-\ell)}\) and \(\widehat q_{d,B}^{(-\ell)}\) are trained outside observational fold \(\ell\), while every experimental observation is evaluated with each of the \(L\) outcome-bridge fits.  For deterministic \(d_B\in\mathcal D\) and \(\alpha_B\in(0,1)\), put
\[
 e_B=\left\lfloor\frac{\alpha_BB}{c_{E,d_B}}\right\rfloor,
 \qquad
 o_B=\left\lfloor\frac{(1-\alpha_B)B}{c_{O,d_B}}\right\rfloor,
 \qquad n_B=e_B+o_B,
\]
and let \(P_{B,d_B,\alpha_B}^{\circ}\) be the artificial two-source law with source probabilities \(e_B/n_B\) and \(o_B/n_B\) and source-conditional laws \(P(\cdot\mid G=E)\) and \(P(\cdot\mid G=O)\).  Write \(\|\cdot\|_{2,B}^{\circ}=\|\cdot\|_{L_2(P_{B,d_B,\alpha_B}^{\circ})}\).

Suppose that along every deterministic sequence satisfying
\[
 |\alpha_B-\alpha_{d_B}^*(P)|\longrightarrow0,
\]
uniformly over \(1\le\ell\le L\), with
\(\Delta h_{B,\ell}=\widehat h_{d_B,B}^{(-\ell)}-h_{d_B}\) and
\(\Delta q_{B,\ell}=\widehat q_{d_B,B}^{(-\ell)}-q_{d_B}\),
\[
 \|T_{d_B}\Delta h_{B,\ell}\|_{2,B}^{\circ}=O_P(\delta_{h,B}),
 \qquad
 \|\Delta h_{B,\ell}\|_{2,B}^{\circ}=O_P(\rho_{h,B}),
\]
\[
 \|T_{d_B}^{\star}\Delta q_{B,\ell}\|_{2,B}^{\circ}=O_P(\delta_{q,B}),
 \qquad
 \|\Delta q_{B,\ell}\|_{2,B}^{\circ}=O_P(\rho_{q,B}),
\]
where all four deterministic rates are \(o(1)\) and
\[
 \min\{\delta_{h,B}\rho_{q,B},\rho_{h,B}\delta_{q,B}\}
 =o(B^{-1/2}).
\]
Also suppose the foldwise multiplier errors obey
\[
\begin{aligned}
 \max_{1\le\ell\le L}\bigl(&|\mathbf 1\{G=O\}\Delta q_{B,\ell}
       (Y-h_{d_B})\|_{2,B}^{\circ}
 +\|\mathbf 1\{G=O\}q_{d_B}\Delta h_{B,\ell}\|_{2,B}^{\circ}\\
 &+\|\mathbf 1\{G=O\}\Delta q_{B,\ell}\Delta h_{B,\ell}\|_{2,B}^{\circ}
 \bigr)=o_P(1),
\end{aligned}
\]
where every bridge in this display is evaluated at its design-specific arguments.  This condition concerns only empirical multiplier stability and does not replace the preceding weak--strong product condition.

For a deterministic design and allocation, define the arm-normalized cross-fitted estimator by
\[
\begin{aligned}
 \widehat\mu_{B,d_B}(a)
 &=\frac1L\sum_{\ell=1}^L
 \frac{\sum_{i\in E_B}\mathbf 1\{A_i=a\}
 \widehat h_{d_B,B}^{(-\ell)}(S_{3d_B,i},S_{2d_B,i},a,X_i)}
 {\sum_{i\in E_B}\mathbf 1\{A_i=a\}}\\
 &\quad+\frac1L\sum_{\ell=1}^L
 \frac{\sum_{i\in O_{B,\ell}}\mathbf 1\{A_i=a\}
 \widehat q_{d_B,B}^{(-\ell)}(S_{2d_B,i},S_{1d_B,i},a,X_i)
 \{Y_i-\widehat h_{d_B,B}^{(-\ell)}(S_{3d_B,i},S_{2d_B,i},a,X_i)\}}
 {\sum_{i\in O_{B,\ell}}\mathbf 1\{A_i=a\}},
\end{aligned}
\]
and set \(\widehat\tau_B(d_B,\alpha_B)=\widehat\mu_{B,d_B}(1)-\widehat\mu_{B,d_B}(0)\).  Give this estimator an arbitrary fixed measurable value whenever an arm denominator is zero.  The implemented \(\widehat\tau_B\) is \(\widehat\tau_B(\widehat d_m,\widehat\alpha_m)\).

Let \(m\to\infty\) and \(B\to\infty\), assume \(\underline v_m=o(1)\), and assume \(\mathcal P_m\) is independent of the main study.  Suppose
\[
 \max_{d\in\mathcal D}\max\{|\widehat V_{E,d}-V_{E,d}|,
 |\widehat V_{O,d}-V_{O,d}|\}=O_P(r_{V,m}),
 \qquad r_{V,m}=o(1),
\]
and suppose
\[
 \max_{d\in\mathcal D,\,g\in\{E,O\}}
 |\widetilde V_{g,d,B}-V_{g,d}|=O_P(r_{\widetilde V,B}),
 \qquad r_{\widetilde V,B}=o(1).
\]
Here \(\widetilde V_{g,d,B}\) is the foldwise empirical second moment of the plug-in version of the exact arm-normalized score \(\phi_{g,d}\), not of the differently normalized score printed in Imbens--Kallus--Mao--Wang Theorem~5.

For the selected design and realized counts, on
\[
 \mathsf G_{m,B}=\{N_{E,B}\ge1,\ N_{O,B}\ge1,\ \widehat s_B^2>0\},
\]
set
\[
 s_{0,B}^2=B\left\{\frac{V_{E,\widehat d_m}}{N_{E,B}}
 +\frac{V_{O,\widehat d_m}}{N_{O,B}}\right\},
 \qquad
 \widehat s_B=(\widehat s_B^2)^{1/2},
\]
using the nonnegative square root, and give every ratio and studentized statistic below an arbitrary fixed measurable value on \(\mathsf G_{m,B}^c\).  Then
\[
 P(\mathsf G_{m,B})\longrightarrow1,
 \qquad
 s_{0,B}^2-\mathcal V_{\widehat d_m}^*(P)=o_P(1),
 \qquad
 \widehat s_B^2-\mathcal V_{\widehat d_m}^*(P)=o_P(1),
\]
and
\[
 \frac{\sqrt B(\widehat\tau_B-\tau)}{\widehat s_B}
 \rightsquigarrow\mathcal N(0,1).
\]
If \(z_{1-\eta/2}\) is the \((1-\eta/2)\)-quantile of the standard normal distribution and
\[
 I_B=\left[
 \widehat\tau_B-z_{1-\eta/2}\frac{\widehat s_B}{\sqrt B},
 \widehat\tau_B+z_{1-\eta/2}\frac{\widehat s_B}{\sqrt B}
 \right],
\]
then \(P\{\tau\in I_B\}\to1-\eta\).